%==============================================================================
%  STAT 5200 - Applied Econometrics I
%  Homework 4
%
%  Compile with:  pdflatex "HW 3 - Report.tex"  (run twice for refs)
%
%  Directory layout assumed by this file:
%     HW 4 - Report.tex
%     R/        (scripts, if Q1 needs descriptive statistics)
%     figures/  (generated figures)
%     output/   (captured console output)
%==============================================================================
\documentclass[11pt]{article}

\usepackage[margin=1in]{geometry}
\usepackage[T1]{fontenc}
\usepackage[utf8]{inputenc}
\usepackage{amsmath,amssymb}
\usepackage{booktabs}
\usepackage{graphicx}
\usepackage{xcolor}
\usepackage{listings}
\usepackage[font=small,labelfont=bf]{caption}
\usepackage[hidelinks]{hyperref}

%------------------------------------------------------------------------------
%  Code presentation
%
%  Code is set as selectable text on a fixed character grid, source and console
%  output are given visually distinct styles, and source lines are kept to 76
%  characters so that nothing wraps.
%------------------------------------------------------------------------------
\definecolor{codekeyword}{RGB}{0,86,148}
\definecolor{codestring} {RGB}{163,21,21}
\definecolor{codecomment}{RGB}{106,115,125}
\definecolor{codebg}     {RGB}{248,248,250}
\definecolor{coderule}   {RGB}{220,220,226}

\lstdefinestyle{Rcode}{
  language        = R,
  basicstyle      = \ttfamily\footnotesize,
  keywordstyle    = \color{codekeyword}\bfseries,
  stringstyle     = \color{codestring},
  commentstyle    = \color{codecomment}\itshape,
  numbers         = left,
  numberstyle     = \ttfamily\scriptsize\color{codecomment},
  numbersep       = 9pt,
  backgroundcolor = \color{codebg},
  frame           = single,
  rulecolor       = \color{coderule},
  framesep        = 4pt,
  xleftmargin     = 15pt,
  xrightmargin    = 0pt,
  columns         = fixed,
  basewidth       = 0.5em,
  keepspaces      = true,
  tabsize         = 2,
  showstringspaces = false,
  breaklines      = true,
  breakatwhitespace = true,
  postbreak       = \mbox{\textcolor{codekeyword}{$\hookrightarrow$}\space},
  captionpos      = b,
  abovecaptionskip = 6pt,
}

\lstdefinestyle{Routput}{
  language        = {},
  basicstyle      = \ttfamily\footnotesize,
  keywordstyle    = \ttfamily,
  commentstyle    = \ttfamily,
  stringstyle     = \ttfamily,
  numbers         = none,
  backgroundcolor = \color{white},
  frame           = leftline,
  rulecolor       = \color{coderule},
  framesep        = 6pt,
  xleftmargin     = 20pt,
  breaklines      = true,
  captionpos      = b,
  abovecaptionskip = 6pt,
}

\lstset{style=Rcode}
\renewcommand{\lstlistingname}{Listing}

% Short inline code fragments.
\newcommand{\code}[1]{\texttt{#1}}

% Operators used throughout.
\DeclareMathOperator{\Var}{Var}
\DeclareMathOperator{\Cov}{Cov}
\DeclareMathOperator{\Avar}{Avar}
\DeclareMathOperator{\sd}{sd}
\DeclareMathOperator*{\plim}{plim}

% Placeholder marker. Every occurrence should be gone before submission;
% search the source for "\todo" to find what is still outstanding.
\newcommand{\todo}[1]{\textcolor{red}{\textbf{[TODO:} #1\textbf{]}}}

%------------------------------------------------------------------------------
\title{\vspace{-2em}Homework 4\\[0.3em]
       \large STAT 5200 --- Applied Econometrics I --- Fall 2026}

\author{Apostolos Kyratzis - 13130833}
\date{October 19, 2026}

\begin{document}
\maketitle

%==============================================================================
\section{Q1: Replication and project proposal}
\label{sec:q1}
%==============================================================================

\todo{no more than one page in total}

\subsection{The problem to be considered}
\todo{what question the project will ask}

\subsection{The data set}
\todo{which data, where it comes from, and confirmation that it is already
on hand; the question asks for descriptive statistics of the variables of
interest to have been computed}

\subsection{The paper to be replicated}
\todo{which paper, and why it helps with the literature and the data}

\clearpage
%==============================================================================
\section{Q2: Paper question --- difference-in-differences}
\label{sec:q2}
%==============================================================================

\subsection{The paper}
\todo{published in a journal from the field, contains
``difference-in-differences'', reports DiD results; confirm the venue before
writing the citation}

\subsection{Citation}
\todo{author, title, year, journal, volume, issue, pages, DOI}

\subsection{The main research question}
\todo{one paragraph}

\subsection{The empirical model}
\todo{identify the outcome, the treatment variable, the controls and fixed
effects in the main specification, and the total number of observations}

\subsection{The source of treatment}
\todo{the event, policy change or intervention}

\subsection{Treatment and control groups}
\todo{who is in each, and why the control group is a reasonable comparison}

\subsection{The coefficient of interest}
\todo{what it measures and how the authors interpret it}

\subsection{The parallel trends assumption}
\todo{what must hold for the estimate to be causal, in this paper's context}

\subsection{Evidence for parallel trends}
\todo{pre-treatment trends, event-study specification, or whatever they show}

\subsection{A story under which parallel trends fails}
\todo{a plausible case where the two groups would have diverged anyway}

\clearpage
%==============================================================================
\section{Q3: Paper question --- synthetic control}
\label{sec:q3}
%==============================================================================

\subsection{The paper}
\todo{published in a journal from the field, uses the synthetic control
method, reports results from it}

\subsection{Citation}
\todo{author, title, year, journal, volume, issue, pages, DOI}

\subsection{The main research question}
\todo{one paragraph}

\subsection{The treated unit and the treatment}
\todo{which unit, and what event or policy constitutes treatment}

\subsection{The donor pool}
\todo{which units it contains}

\subsection{How the synthetic control is constructed}
\todo{the variables and pre-treatment outcomes used to choose the weights}

\subsection{The estimated treatment effect}
\todo{how it follows from comparing the treated unit with its synthetic
control after treatment}

\subsection{Pre-treatment fit}
\todo{how closely the synthetic control tracks the treated unit beforehand,
and what the authors show}

\subsection{Placebo and permutation tests}
\todo{what they do and what they establish}

\subsection{A criticism of the design}
\todo{why the donor pool may not give a good counterfactual, or what else
happened around the treatment date}

\clearpage
%==============================================================================
\section{Q4: Paper question --- regression discontinuity}
\label{sec:q4}
%==============================================================================

\subsection{The paper}
\todo{published in a journal from the field, contains ``regression
discontinuity'', reports RDD results}

\subsection{Citation}
\todo{author, title, year, journal, volume, issue, pages, DOI}

\subsection{The main research question}
\todo{one paragraph}

\subsection{The empirical model}
\todo{identify the outcome, the running variable, the cutoff, any controls,
and the total number of observations}

\subsection{How the cutoff determines treatment}
\todo{why crossing it changes treatment status or its probability}

\subsection{Sharp or fuzzy}
\todo{which, and why; if fuzzy, how the cutoff shifts the probability of
treatment}

\subsection{The coefficient of interest}
\todo{what it measures and how the authors interpret it}

\subsection{Bandwidth and robustness}
\todo{the bandwidth used, and whether results survive other bandwidths and
functional forms}

\subsection{The identifying assumption}
\todo{why units just either side of the cutoff are comparable}

\subsection{Evidence for validity}
\todo{manipulation tests, the density around the cutoff, continuity of
predetermined characteristics, graphical analysis}

\subsection{A criticism of the design}
\todo{manipulation of the running variable, another policy at the same
cutoff, or another violation of continuity}

\end{document}
